\section{Inside an Epoch}
\label{sec:epoch}
Within an epoch RAI votes like Nano or FastPay, with two twists that pay off at the boundary: a validator never switches sides, and a final vote is preceded by a notarization.

\para{Confirming a block.}
Alice signs block $A$ at position 7 of her account and sends it to the six validators. Each validator checks the block against Alice's chain as it knows it (parent, balance, no duplicate receive) and, if it has not voted at position 7 in this epoch, casts a first vote for $A$ (\ref{r:one}). With four first votes $A$ is \emph{notarized}; a validator that holds the notarization and first-voted $A$ casts its final vote. Four final votes make $A$ final. If five first votes arrive from validators that have installed the previous checkpoint, $A$ is final at once, in one round trip; this is the common case, and the reason RAI's median latency in Section~\ref{sec:eval} is below Nano's, which always needs a second round.

\para{Why the thresholds are safe.}
Any two groups of four validators out of six share at least two, and at most one of those is Byzantine. A correct validator supports one block per position (\ref{r:one}), so two conflicting blocks cannot both reach four supporters. Four supporters for $A$ therefore \emph{exclude} every rival: no rival can be notarized or final in this epoch. We call a set of $q$ supporters an \emph{exclusion witness}; it is the one piece of evidence the boundary rules will ask for again and again. Five first votes for $A$ leave at most one correct validator free to support a rival, so a fast-final block has no rival with three honest supporters either; this is what makes the recovery threshold $r=3$ of~\ref{r:build} work.

\para{A double spend.}
Suppose Alice also signs $A'$ at position 7, paying someone else, and shows $A$ to three validators and $A'$ to the other three (Figure~\ref{fig:fork}). Neither reaches four. In Nano this position is stuck: final votes cannot be withdrawn, no block at position 7 will ever confirm, and every later block of Alice's account waits behind it. In RAI the next checkpoint decides what may still happen at position 7 (\ref{r:build}). It is built from five of the six reports, which show three first votes for one side and two for the other, say three for $A$. $A$ is kept under a \emph{recovery lock} of origin $e$ and $A'$ is dropped: from now on no validator admits $A'$ or any other rival of $A$ at position 7, so $A'$ can never become final. Had the reports shown three against three, the block with the smaller hash would be kept.

\para{How the fork ends.}
The checkpoint decides the position (\ref{r:install}). A validator that has installed it votes only $A$ at position 7, and a validator that has seen both $A$ and $A'$ casts no vote there before it has the checkpoint (\ref{r:cross}). So $A$ collects four first and four final votes and becomes final by the second epoch of normal operation after the fork, whatever Alice or the Byzantine validator does (Theorem~\ref{thm:live}), and Alice's account continues from it. She need not act; to pay sooner she can sign a block $C$ on top of $A$, admitted once $A$ has four first votes again in the new epoch, and finalizing $C$ finalizes $A$ (Figure~\ref{fig:fork}). A block on top of $A'$ is refused.

\para{Why the checkpoint must decide.}
Suppose validators could still vote $A'$ whenever they held four epoch-$e$ supporters for it, the witness that releases $A$'s lock; that takes a Byzantine validator that supported both sides. Then in every epoch some correct validators vote $A$ and some $A'$, the Byzantine validator withholds the one vote either side needs, and the position stays open forever, as in Nano (supplement). With the checkpoint deciding, the witness still lets a validator accept a certificate for $A'$ that already exists, so nothing final is contradicted, but it earns $A'$ no new vote.

\begin{figure}[t]
\centering
\begin{tikzpicture}[>=Stealth,font=\footnotesize]
\node[final] (a6) at (0,0) {$A_6$};
\node[blk] (a) at (1.2,0.55) {$A$};
\node[blk] (ap) at (1.2,-0.55) {$A'$};
\draw[->] (a) -- (a6); \draw[->] (ap) -- (a6);
\node[val,anchor=west] at (1.65,0.55) {first votes $V_1,V_2,V_3$};
\node[val,anchor=west] at (1.65,-0.55) {first votes $V_4,V_5,V_6$};
\node[val,anchor=north] at (0,-0.3) {final};
\draw[->,thick,cPause] (4.1,0) -- (4.9,0) node[midway,above,val] {close $e$};
\node[locked] (ka) at (5.4,0.55) {$A$};
\node[val,anchor=west] at (5.75,0.55) {lock, origin $e$};
\node[val,anchor=west,text=cBad] at (5.75,-0.55) {$A'$ dropped};
\node[final] (fa) at (5.4,-1.3) {$A$};
\node[final] (fc) at (6.4,-1.3) {$C$};
\draw[->] (fc) -- (fa);
\draw[->,thick,cFinal] (ka) -- node[val,left]{owner} (fa);
\node[val,anchor=west] at (6.75,-1.3) {Alice adds $C$};
\end{tikzpicture}
\caption{A double spend at position 7. Neither block reaches four first votes. The five selected reports show three votes for $A$ and two for $A'$: the checkpoint locks $A$ for recovery and drops $A'$, which can never become final. In the next epochs every validator votes $A$, which becomes final; Alice may also add $C$ on top of $A$.}
\label{fig:fork}
\end{figure}

\para{Durable votes.}
A correct validator writes its first vote and its final vote to disk before releasing them, so that a crash and restart cannot make it vote twice. Our prototype syncs these writes; Section~\ref{sec:eval} reports what that costs.
